\section{Introduction}
\label{sec:introduction}

Animal Re-identification (ReID) aims to recognize and match a specific individual across different non-overlapping camera views, despite variations in pose or viewpoint.
It is an essential component for automated tracking in wildlife conservation and livestock management.

However, in real-world deployments, camera placement is often constrained by physical limitations in closed spaces or wild environments.
This is particularly challenging in Aerial-Ground ReID (AG-ReID) scenarios, where models must associate individuals across drastic perspective changes.
Identifying which viewpoints maximize identification accuracy is therefore crucial for optimizing sensor placement configurations.

While viewpoint dependencies have been extensively studied in human ReID~\cite{sunDissectingPersonReIdentification2019a}, equivalent insights for animals remain lacking.
This gap is detrimental, as animals often exhibit asymmetric coat patterns, making their identification highly sensitive to viewpoint shifts.
Existing datasets lack precise angular annotations and cannot perform systematic analysis of viewpoint impact.

To bridge this gap, we introduce the \textbf{M}ulti-view \textbf{O}riented \textbf{O}bservation (MOO), a large-scale synthetic AG-ReID dataset.
MOO contains 1,000 synthetic cattle identities rendered from 128 uniformly sampled viewpoints spanning $360^\circ$ in azimuth and $-25^\circ$ to $90^\circ$ elevation (relative to the horizontal side view).
By isolating identity through asymmetric patterns, we systematically analyze viewpoint sensitivity.
Ultimately, this work lays the foundation for future model development on geometric robustness in animal ReID.

Our contributions are threefold:
\begin{itemize}
    \item We present the MOO dataset, a large-scale synthetic benchmark for patterned animal ReID with precise azimuth and elevation annotations.
    \item We provide the first systematic quantification of elevation impact on patterned animal ReID, identifying a $30^\circ$ elevation angle as a key threshold for generalization.
    \item We demonstrate that MOO serves as a robust pre-training source, yielding consistent performance gains on real-world datasets in both zero-shot and supervised scenarios.
\end{itemize}

\begin{figure*}[tbp]
    \centering
    \includegraphics[width=1.0\linewidth]{gabstract}
    \caption{MOO: A synthetic dataset for systematic analysis and real-world transferability. The workflow encompasses (1) individual generation; (2) dataset rendering; and (3) systematic analysis and real-world transfer.}
    \label{fig:gabstract}
\end{figure*}
